%!TEX encoding = UTF8
%!TEX root = 0-notes.tex

\chapter{Fractions rationnelles}
\label{chap:fonctions-rationnelles}

\section{Corps des fractions}

\nt{
	L'anneau $\Z$ n'est pas un corps car il n'est pas possible de diviser à l'interieur de celui-ci.
	Quel est le plus petit sur-ensemble de $\Z$ dans lequel il est possible de diviser ?
	La réponse est bien sûr $\Q$, ensemble dans lequel les éléments sont notés $\frac{a}{b}$ avec $a\in\Z$ et $b\in\Z^*$.
	
	Bien sûr, $\frac{a}b = ab^{-1}$, donc l'ajout d'un inverse multiplicatif de chaque élément non nul de $\Z$ est nécessaire et suffisant à créer $\Q$.
	De plus, cette notation fractionnelle nécessite que l'anneau $\Z$ soit commutatif, sans quoi nous ne pourrions pas savoir si $b$ divise $a$ à gauche ou à droite.
	
	Notons en outre qu'il n'est pas souhaitable que $b$ divise zéro, auquel cas, si $bx = 0$ avec $x\neq0$, alors $\frac{a}{b} \times \frac{1}x = \frac{a}0$.
	Ceci n'arrive bien sûr pas dans $\Z$ qui est un \emphindex{anneau intègre}.
}

\exe{}{
	Montrer que $\Q$ est le plus petit corps commutatif, pour l'inclusion et à isomorphisme près, contenant $\Z$.
}{exe:Q-smallest-field}{
	Notons $F$ un corps commutatif contenant $\Z$.
	L'application $\iota : \frac{a}{b} \mapsto ab^{-1}$ est une injection de $\Q$ dans $F$.
	Sa bonne définition est due au fait que $\iota\left(\frac{ac}{bc}\right) = (ac)(bc)^{-1} = a(cc^{-1})b^{-1} = \iota\left(\frac{a}{b}\right)$.
}

\dfn{Corps des fractions}{
	Soit $R$ un anneau intègre et commutatif.
	Alors le \emphindex{corps des fractions} de $R$ est donné par
		\[ \Frac(R) = \Bigset{ \frac{a}{b} \tq a \in R \et b \in R^* }. \]
}{dfn:corps-frac}

\dfn{}{
	Le corps des fractions de l'anneau des polynômes $\K[X]$ est noté
		\[ K(x) = \left\{ \frac{P(x)}{Q(x)} \tq P, Q \in\K[X], Q\neq 0 \right\}. \]
	Un élément est appelé \emphindex{fraction rationnelle}.
}{dfn:frac-rationnelle}

\section{Décomposition en éléments simples}

\subsection{Cas général}

\dfn{élément simple}{
	Soit $p\in\K[X]$ un polynôme irréductible.
	Alors un élément simple associé à $p$ est une fraction rationnelle qui prend la forme
		\[ \frac{f}{p^n}, \]
	avec $n\in\N^*$ un entier naturel, et $f\in\K[X]$ un polynôme vérifiant $\deg(f) < \deg(p)$.
}{dfn:el-simple}

\thm{}{
	Soit $\frac{P}Q\in\K(x)$ une fraction rationnelle.
	Supposons que $Q$ se scinde en produit d'irréductibles comme
		\[ Q = p_1^{n_1} p_2^{n_2} \cdots p_k^{n_k}. \]
	Alors, la fraction $\frac{P}Q$ se décompose en une partie entière $E\in\K[X]$, et en une somme d'éléments simples de la forme $\frac{f}{p^n}$ où $n$ est inférieur ou égal à la puissance de $p$ dans la décomposition de $Q$.
	
	Autrement dit, nous pouvons écrire,
		\[ \frac{P}{Q} = \underbrace{E}_{\text{partie entière dans $\K[X]$}} + \sum_{i=1}^{k} \sum_{n=1}^{n_k} \underbrace{\frac{f_{in}}{p_i^n}}_{\text{élément simple}}, \]
	où $\deg(f_{in}) < \deg(p_i)$ pour chaque $i, n$.
}{thm:el-simples-K}

\subsection{Le cas de $\R[x]$}

Sur $\R$, les polynômes irréductibles sont de degrés exactement 1 ou 2, d'après le corollaire \ref{cor:decomp-Rx} :
	\[ Q(x) = \prod_{u, n} (x-u)^n \prod_{v, w, n} (x^2 + vx + w)^n. \]
Par conséquent, les facteurs simples prennent la forme
	\begin{align*}
		\frac{a}{(x-u)^n} && \ou && \frac{ax+b}{(x^2 + vx+ w)^n},
	\end{align*}
où $a, b, u, v, w\in\R$ sont des constantes réelles.



\subsection{Le cas de $\C[x]$}

Sur $\C$, les polynômes irréductibles sont de degrés exactement 1, d'après le corollaire \ref{cor:decomp-C} :
	\[ Q(x) = \prod_{u, n} (x-u)^n. \]
Par conséquent, les facteurs simples prennent la forme
	\begin{align*}
		\frac{a}{(x-u)^n}
	\end{align*}
où $a, u\in\C$ sont des constantes complexes.

\section*{Exercices}
\addcontentsline{toc}{section}{Exercices}


\exe{}{
	Considérons la {suite de Fibonacci} $(F_n)_{n \geq 0}$ définie récursivement par 
		\begin{align*} F_0 = 0, && F_1 = 1, && F_{n+2} = F_{n+1} + F_{n} \quad (n \geq 0). \end{align*}
	Posons $F$ la fonction génératrice suivante, où $x\in\R$.
		\[ F(x) = \sum_{n=0}^{\infty} F_n x^n = F_0 + F_1 x + F_2 x^2 + F_3 x^3 + \dots \]
	%Nous admettons que $F$ converge absolument pour $x$ suffisamment petit.
	Nous étudions d'abord la suite de Fibonacci et la convergence de la série $F$.
	\begin{enumerate}
		\item
		Écrire les quatre premiers termes de la suite de Fibonacci.
		\item
		Combien d'opérations arithmétiques sur les nombres réels (addition, soustraction, multiplication, division) sont nécessaires pour calculer $F_{1024}$ ?
		\item
		Montrer que $F_n \leq 2^n$ pour tout $n\in\N$.
		\item
		En déduire un rayon $R>0$ tel que la série $F(x)$ converge absolument pour tout $|x|<R$.
	\end{enumerate}
	Les manipulations de sommes usuelles sont permises par la convergence absolue ($|x|<R$).
	\begin{enumerate}[resume]
		\item
		Montrer que
			\[ x^2 F(x) + x F(x) = F(x) - x \]
		et en déduire que 
			$F(x) = \frac{x}{1-x-x^2}.$
		\item
		Trouver les nombres complexes $z_1, z_2, w\in\C$ tels que
			\[ F(x) = w \left( \frac1{1-z_1 x} - \frac1{1-z_2x} \right). \]
	%	pour tout $|x|<R$.
	\end{enumerate}
	La décomposition en éléments simples d'une fonction génératrice permet d'obtenir une expression fonctionnelle de ses coefficients grâce aux résultats classiques suivants.
	\begin{enumerate}[resume]
		\item
		Démontrer que, pour $|t|<1$,
			\[ \frac{1}{1-t} = \sum_{n=0}^{\infty}t^n. \]
		\item
		Montrer que l'égalité des séries absolument convergentes de rayon de convergence non nul $\sum_{n=0}^{\infty} a_k x^k = \sum_{n=0}^{\infty} b_k x^k$ implique $a_n = b_n$ pour tout $n\in\N$.
	\end{enumerate}
	Nous concluons à l'aide de la décomposition en éléments simples de $F(x)$.
	\begin{enumerate}[resume]
		\item
		Montrer, à l'aide des questions précédentes, que le $n$-ième nombre de Fibonacci est égal à
		\begin{align*}
			F_n = \frac{\phi^{n} - (\overline{\phi})^{n}}{\sqrt5} && \text{où } \phi = \frac{1+\sqrt5}2 \et \overline{\phi} = \frac{1-\sqrt5}2.
		\end{align*}
		\item
		En déduire que $F_n$ est l'entier le plus proche de $\frac{\phi^n}{\sqrt5}$ pour tout $n$ suffisamment grand.
		\item
		Combien d'opérations arithmétiques sur les nombres réels (addition, soustraction, multiplication, division) sont nécessaires pour calculer $F_{1024}$ ?
	\end{enumerate}
}{exe:Fibo}{
	\begin{enumerate}
		\item
		Nous avons $F(x) = F_0 + F_1x + F_2x^2 + F_3x^3 + \dots = x + x^2 + 2x^3 + \dots$.
	
		\item
		Le membre de gauche est égal à
		\begin{align*}
			x^2 F(x) + x F(x) &= \sum_{n=0}^{\infty} F_n x^{n+2} + \sum_{m=0}^{\infty} F_mx^{m+1} \\
				&= \sum_{n=0}^{\infty} F_n x^{n+2} + F_0 x + \sum_{m=0}^{\infty} F_{m+1} x^{m+2} \\
				&= \sum_{n=0}^{\infty} (F_n + F_{n+1}) x^{n+2} \\
				&= \sum_{n=0}^{\infty} F_{n+2}x^{n+2} \\
				&= \sum_{n=0}^{\infty} F_n x^n - F_0 - F_1x \\
				&= F(x) - x
		\end{align*}
		En regroupant les termes en $F(x)$, nous obtenons $F(x) (1 - x - x^2) = x$.
		Diviser par $1-x-x^2$ conclut.
		
		\item
		Le membre de droite est égal à
		\begin{align*}
			 \frac{1}{\sqrt5} \left( \frac1{1-\phi x} - \frac1{1-\overline{\phi}x} \right) &= 
			 \frac{1}{\sqrt5} \frac{1-\overline{\phi}x - (1-\phi x)}{(1-\phi x)(1-\overline{\phi}x)}  \\
			 &= \frac{1}{\sqrt5} \frac{(\phi-\overline{\phi})x}{1 - (\phi + \overline{\phi})x + \phi\overline{\phi} x^2}  \\
			 &= \frac{1}{\sqrt5} \frac{\sqrt5x}{1 -x - x^2}  \\
			 &= \frac{1}{1-x-x^2} \\
			 &= F(x)
		\end{align*}
		
		\item
		D'après la question précédente,
		\begin{align*}
			F(x) &= \frac{1}{\sqrt5} \left( \frac1{1-\phi x} - \frac1{1-\overline{\phi}x} \right) \\
			&= \frac{1}{\sqrt5} \left( \sum_{n=0}^{\infty} (\phi x)^n - (\overline{\phi} x)^n \right) \\
			&= \frac{1}{\sqrt5} \sum_{n=0}^{\infty} \left[\phi^n - \overline{\phi}^n \right] x^n \\
			&= \sum_{n=0}^{\infty} \frac{\phi^n - \overline{\phi}^n}{\sqrt5} x^{n}
		\end{align*}
		d'où $F_{n} = \frac{\phi^n - \overline{\phi}^n}{\sqrt5}$.
	
		\item
		Comme $|\overline{\phi}| < 1$ par inspection, le terme $\overline{\phi}^n$ tend vers zéro lorsque $n$ tend vers l'infini.
		Par conséquent, $F_n \approx \phi^n / \sqrt5$.
	\end{enumerate}
}

\exe{}{
	Justifier simplement qu'une suite $(G_n)_{n\geq0}$ définie récursivement par $aG_{n+2} = bG_{n+1} + cG_{n}$,
	où $a, b, c\in\R$ sont trois nombres réels, prend une des deux formes suivantes.
	\begin{enumerate}
		\item 
		Si $b^2 - 4ac \neq 0$, alors
			\[ G_n = \alpha z_1^n + \beta z_2^n, \]
		où $\alpha, \beta\in\C$ et $z_1, z_2\in\C$ sont les racines du polynôme $az^2 + bz + c$.
		\item
		Si $b^2 - 4ac = 0$, alors
			\[ G_n = \alpha z^n + \beta n z^n, \]
		où $\alpha, \beta\in\C$ et $z\in\C$ est l'unique racine du polynôme $az^2 + bz + c$.
	\end{enumerate}
}{exe:Fibo-gen}{

}


\shipoutExercise

%\section*{Solutions}
%\addcontentsline{toc}{section}{Solutions}
%
%\shipoutAnswer



